% 本文件由 LaTeX 排版 Agent 根据有效 translation.json 自动生成。
% author=Gregory of Nyssa; work_id=2911; title=书信集

\chapter{书信集}

% structure: p0001 source=h1 target=omitted reason=page-title-already-rendered-by-parent

第一封 致\Person{优西比乌}\newline{}第二封 致\Place{塞巴斯特}城\newline{}第三封 致\Person{阿布拉比乌}\newline{}第四封 致\Person{辛尼吉乌}\newline{}第五封 推荐信\newline{}第六封 致\Person{斯塔吉里乌}\newline{}第七封 致一位朋友\newline{}第八封 致一位古典学学生\newline{}第九封 邀请函\newline{}第十封 致\Person{利巴尼乌}\newline{}第十一封 致\Person{利巴尼乌}\newline{}第十二封 论其驳斥\Person{欧诺米}之作\newline{}第十三封 致\Place{尼科美底亚}教会\newline{}第十四封 致\Place{梅利特内}主教\newline{}第十五封 致律师\Person{阿德尔斐乌}\newline{}第十六封 致\Person{安斐洛奇乌}\newline{}第十七封 致\Person{尤斯塔提亚}、\Person{安布罗西亚}与\Person{巴西利萨}\newline{}第十八封 致\Person{弗拉维安}\par

\SourceRecord{Translated by H.C. Ogle, H.A. Wilson and W. Moore.；Nicene and Post-Nicene Fathers, Second Series；Vol. 5.；Edited by Philip Schaff and Henry Wace.；Buffalo, NY: Christian Literature Publishing Co.,；1893.；<http://www.newadvent.org/fathers/2911.htm>.}

% structure: p0001 source=h1 target=section reason=page-title

\section{第一封信}

\Emph{\Emph{\Person{圣额我略·尼撒}著}}\par

\Emph{致\Person{欧瑟伯}书}\par

当冬季白日开始延长，太阳升至其行程的上半部时，我们庆祝那透过血肉之幕、将灿烂光芒投射在人类生命上的真光——天主的显现庆节。但现在，当那发光体在半空运行，使昼夜等长时，人类本性从死亡返归生命的上升，便成了这伟大而普世的庆节的主题，凡接受复活奥迹者的一生都联合庆祝之。这个在我信函中提出的主题有何含义？因为在这些普遍的庆节中，按惯例我们以各种方式表达心中深藏的情感，而有些人，如你所知，以各自的礼物证明善意。我们认为理应不让你缺少我们的敬意，而应将我们贫穷的菲薄献礼呈于你那崇高而高贵的灵魂。现在，在这封信中献予你接纳的礼物，便是这封信本身。信中没有一句辞藻华丽或修饰优美的语句，使之在文坛上被视为礼物；但其中包裹着基督徒信德的神秘黄金，如同包裹一般，必须藉这些文字尽可能展开，显出其隐藏的光辉，作为我的赠礼。因此，我们必须回到序曲。为何只有当黑夜达到极致，无可再增时，那掌握宇宙、以己力统御宇宙、甚至不为一切理智之物所包容、却包容万有的那一位，竟进入血肉帐幕的狭窄居所，而他的大能如此配合他的慈爱旨意，甚至如影子般随意志动向显现，以致在创世时大能并不弱于意志，当他渴望俯就我们必死本性的卑下时，也并不缺乏达到此目的的能力，反而确实进入那种状态，却并未使宇宙失去引导？既然这两个季节都需解释——为何他在冬天显现于肉身，却在昼夜等长时，使因罪返回尘土的人重获生命——我将用几句话尽可能解释原因，使我的信成为献于你的礼物。你的明智（当然如此）是否已猜出这些巧合所暗示的奥迹？即黑夜的进展被光的增加所阻止，黑暗的时期开始缩短，白日的长度因逐次增加而延长？因为这一点对未入教者也足够清楚：罪近似黑暗；事实上，圣经如此称呼恶。因此，我们虔诚奥迹开始的季节，便是天主为我们的灵魂所施救恩计划的一种展示。因为正当合宜，当邪恶无度蔓延时，\BibleQuote{在这邪恶的黑夜中，正义的太阳应当升起}，并且在我们这些\BibleQuote{先前在黑暗中行走}的人中，那从将光置于我们心中的那一位所领受的白日，应日益增长；以致在光明中的生命尽可能延长，不断因善的增加而扩展；而在恶中的生命则通过逐渐减少而缩至最小；因为善的增长等同于恶的减少。但发生在昼夜等长时的复活庆节，本身便给了我们这种巧合的解释：即我们不再与恶势均力敌地战斗，恶习与德行不分胜负地角力，而是光明的生命将得胜，偶像崇拜的阴翳随着白日增强而消融。因此，在月亮运行十四天后，复活节展示她正对太阳的光芒，充满了他光辉的一切富足，不容许任何黑暗的间隔在其间出现：因为她在日落后取代其位置，但自己却不在与太阳真光融合之前落下，以致一道光明持续存留，贯穿地球昼夜运行的全部空间，毫无黑暗的阻隔。亲爱的，这讨论是我们从贫乏之手献上的礼物；愿你的整个人生成为不断的庆节和盛日，永不被一丝夜间的黑暗所玷污。\par

\SourceRecord{Translated by H.C. Ogle.；Nicene and Post-Nicene Fathers, Second Series；Vol. 5.；Edited by Philip Schaff and Henry Wace.；Buffalo, NY: Christian Literature Publishing Co.,；1893.；<http://www.newadvent.org/fathers/291101.htm>.}

% structure: p0001 source=h1 target=section reason=page-title

\section{书信第二}

\Emph{\Person{圣额我略·尼森}}\par

\Emph{致\Place{塞巴斯特亚}城}\par

有些与我们同心同德的弟兄告诉我们，那些仇恨和平、暗中毁谤邻舍、不惧怕那伟大而可畏之审判台的人——祂曾宣告，在我们都必须面对的生命审判中，连虚浮的话也要交账——正传播损害我们的诽谤。他们说，对我们流传的指控如下：我们持守与在\Place{尼西亚}阐明正确纯全信德者相反的意见，并且未经仔细辨别和查问，就接纳了那些先前在\Place{安卡拉}以\Person{玛尔塞鲁}之名聚集的人进入天主教会的共融。因此，为使谎言不致压倒真理，我们在另一封信中已对加诸我们的指控作了充分的辩护，并在主前声明：我们既没有偏离圣教父们的信德，也没有在处理从\Person{玛尔塞鲁}的共融转至教会的共融那些人时未经仔细辨别和查问；我们所做的一切，都是在东方的正统派和牧职中的弟兄们将对这些人的审议托付给我们，并认可我们的行动之后才进行的。然而，自从我们写了那篇书面辩护之后，又有一些与我们同心的弟兄恳请我们亲口单独宣示我们所深信不疑的信仰，我们既遵循圣神的默感和教父的传统，也就认为有必要简要论述这些要点。我们承认，主教导祂门徒的教义——当祂将虔敬的奥迹交付给他们时——是正确纯全信德的基础和根源；我们也不相信有任何比这传统更高超或更安全的道理。主的教义是这样的：祂说：\BibleQuote{你们要去，教导万民，因父、及子、及圣神之名给他们授洗。}既然对于那些从死亡重生而获得永生的人，正是通过天主圣三，将赋予生命的能力赐予那些因信德配享恩宠的人；同样，如果在拯救性的圣洗中遗漏了天主圣三任何一位的名字，恩宠便不完全——因为重生的圣事不能没有圣神而仅由圣子和圣父完成；也不能压抑圣子的名字而在圣父和圣神中赋予生命的完美恩赐；也不能排除圣神而在圣父和圣子中实现复活的恩宠。为此，我们将我们的一切希望和灵魂得救的确信，寄托在这三个被这些名字所认识的位格上；我们相信我们的主耶稣基督的父，祂是生命的泉源；相信父的独生子，祂是生命的创始者，如宗徒所说；并相信天主的圣神，主曾论及祂说：\BibleQuote{使人生活的是圣神。}并且，正如我们所说，由于我们这些从死亡中被赎的人，通过信仰圣父、圣子和圣神获得了不朽的恩宠，我们受此引导而相信：没有任何奴仆性的、被造的、或与圣父的尊威不相称的事物，可以与天主圣三在思想上相关联；因为，我说，我们的生命是通过信仰天主圣三而临于我们的，它源自万有的天主，经过圣子，并由圣神在我们内运作。既然如此确信，我们便按所受的命令受洗，我们按所接受洗礼的方式相信，我们按所相信的持守；因此，我们的洗礼、我们的信德和我们的赞美颂，一致归于父、子及圣神。但若有人提及两个或三个神，或三个神性，愿他受诅咒。若有人追随\Person{阿里乌}的谬误，说圣子或圣神从无中产生，愿他受诅咒。但凡按照真理的准则行事，承认三个位格，并虔敬地按照各自的特性认识它们，相信只有一个神性、一个美善、一个统治、一个权能和能力，既不抹煞独一至高统治的超越性，也不堕入多神论，不混淆位格，不以异质不同的元素构成天主圣三，而是单纯地接受信德的教义，将全部得救的希望寄托于父、子及圣神——这样的人，按我们的判断，是与我们同心同德的，我们也相信将与他们在主内同得基业。\par

\SourceRecord{Translated by H.C. Ogle.；Nicene and Post-Nicene Fathers, Second Series；Vol. 5.；Edited by Philip Schaff and Henry Wace.；Buffalo, NY: Christian Literature Publishing Co.,；1893.；<http://www.newadvent.org/fathers/291102.htm>.}

% structure: p0001 source=h1 target=section reason=page-title

\section{第三封信}

\Emph{圣额我略·尼撒 著}\par

\Emph{致\Person{阿布拉比乌}}\par

主按祂的适当与正义，因你们的祈祷陪伴我们，使我们安然渡过，我将告诉你们一个祂慈爱的明显标记。因为当太阳正照在我们离开\Place{厄尔苏}后方的地点时，忽然乌云密布，晴空变为深沉的黑暗。接着，一股冷风穿过云层，带来细雨，带着十分潮湿的感觉吹在我们身上，预示着前所未见的大雨；左边雷声不断，闪电与雷声交替，一阵雷声之后，下一阵之前，前后左右的山峦都被云笼罩。已经有沉重的乌云悬在我们头顶，被强风吹动，饱含雨水，然而我们，如同古时以色列人在\Place{红海}的神奇渡越，尽管四周被雨包围，却未沾湿而抵达\Place{韦斯特纳}。当我们已在那里找到遮蔽，骡子也得到休息时，天主便向空气发出倾盆大雨的信号。我们在那里停留了大约三四个小时，得到充分休息后，天主再次止住了大雨，我们的车比先前更轻快地前行，因为车轮在表面湿润的泥地上轻易滑过。从那里到我们小城的道路全程沿着河边，顺流而下，河岸村庄连绵不断，都紧邻道路，彼此相距甚近。由于这不间断的居住带，整条路上挤满了人，有的来迎接我们，有的护送我们，在喜悦中混杂着大量的眼泪。此时有微雨，并不令人不快，只够湿润空气；但就在我们到家之前不远，悬在我们头顶的云凝聚成更猛的阵雨，因此我们进入时相当安静，因为没有人预先知道我们的到来。但当我们进入门廊时，车轮在干硬地面上的声音被听到，人们成群出现，仿佛出于某种机械装置，我不知道从哪里、如何而来，如此紧密地拥在我们周围，以至于难以从车上下来，因为没有一尺空地。但当我们艰难地说服他们让我们下车、让骡子通过时，我们被四周拥挤的人群压得喘不过气，以至于他们过分的友善几乎使我们晕厥。当我们靠近门廊内部时，我们看到一条火流涌入教堂；因为贞女的歌咏团手拿蜡炬，正列队沿着教堂入口行进，用火焰点燃了整个地方的辉煌。当我进入里面，与我的百姓一同欢喜和哭泣——因为我目睹人群中的两种情感，自己也经历了这两种情绪——我完成祈祷后，便尽快给您的圣座写下这封信，因极度口渴所迫，这样写完后我便可以照顾身体的需求。\par

\SourceRecord{Translated by H.C. Ogle.；Nicene and Post-Nicene Fathers, Second Series；Vol. 5.；Edited by Philip Schaff and Henry Wace.；Buffalo, NY: Christian Literature Publishing Co.,；1893.；<http://www.newadvent.org/fathers/291103.htm>.}

% structure: p0001 source=h1 target=section reason=page-title

\section{第四封信}

\Signature{圣额我略·尼撒}\par

\Emph{致齐尼吉乌斯}\par

我们有一条诫命，要我们\Quote{与喜乐的人同乐，与哀哭的人同哭}；但这两条诫命中，往往似乎只能实践一条。因为世上\Quote{喜乐的人}极为稀少，不易找到我们能与之分享福乐的人；而处境相反的人却很多。我写这些话作为前言，是因为某种恶毒的力量在长久以来的贵族家庭中上演了一幕悲惨的悲剧。一位名叫西内西乌斯的青年，出身良家，与我自己也有关系，正值青春年华，几乎尚未开始生活，却身陷巨大的危险之中，只有天主有能力拯救他，其次便是您，您掌管着一切生死问题的判决。发生了一件非自愿的不幸——确实，有什么不幸是自愿的呢？而那些对他提出诉讼、要求处以死刑的人，现在已将他的不幸变成了指控的材料。然而，我将尝试通过私人信件来缓和他们的愤怒，使他们倾向怜悯；但我恳求您的仁慈与正义和我们站在一起，使您的善心能胜过那青年的悲惨处境，想尽一切办法，使这年轻人借着您的帮助，战胜攻击他的恶毒力量，脱离危险。我已简要地说了我所要说的；但若详细陈述——为了使我的努力得以成功——那将是我所不该说、也不该由您听我说的话。\par

\SourceRecord{Translated by H.C. Ogle.；Nicene and Post-Nicene Fathers, Second Series；Vol. 5.；Edited by Philip Schaff and Henry Wace.；Buffalo, NY: Christian Literature Publishing Co.,；1893.；<http://www.newadvent.org/fathers/291104.htm>.}

% structure: p0001 source=h1 target=section reason=page-title

\section{第五封信}

\Signature{圣额我略·尼撒}\par

\Emph{推荐书}\par

\Place{马其顿人}的君王在明理之人中最受称颂之处——因为人们称赞他，与其说是由于他战胜\Place{波斯人}和\Place{印度人}的著名战绩，以及他深入到大\Place{大洋}，不如说是由于他说过他的财宝在于他的朋友——就这方面而言，我敢将自己与他神奇的事迹相比，我也有充分理由说出这样的感慨。如今，由于我富有友谊，也许在那种财富上我甚至超过了那位以此自夸的伟人。因为有谁像你之于我那样是你的一位朋友呢？你不断努力在每一种德行上超越自己。因为确实，若有人看到我的年纪和你的生平，绝不会指责我在此话中阿谀奉承：因为白发苍苍不宜阿谀，老年不适于奉承；至于你，即使你始终适宜受奉承，然而赞美也不会落入奉承的嫌疑，因为你的生活在言语之前就已经彰显出对你的赞美。但是，既然人在福泽丰盈时，懂得如何使用自己所拥有的是特别的恩赐，而善用富余的最好方式就是让朋友们分享；又因为我的爱子\Person{亚历山大}是一位与我以完全真诚相系的朋友，所以请听劝，向他展示我的财宝，并不仅展示给他，而且让他随意丰厚地享用，通过在他前来求你担任其保护人的那些事上，向他伸出你的保护之手。他将亲口告诉你一切。因为这样比我写信详细叙述更好。\par

\SourceRecord{Translated by H.C. Ogle.；Nicene and Post-Nicene Fathers, Second Series；Vol. 5.；Edited by Philip Schaff and Henry Wace.；Buffalo, NY: Christian Literature Publishing Co.,；1893.；<http://www.newadvent.org/fathers/291105.htm>.}

% structure: p0001 source=h1 target=section reason=page-title

\section{第六书}

\Emph{\Person{圣尼撒的额我略}著}\par

\Emph{致\Person{斯塔吉里乌斯}}\par

有人说，剧院的变戏法者能制造出这样的一种奇迹，我现在就来描述一下。他们取一段历史叙述或一个古老故事作为戏法的底本，然后用动作向观众讲述这个故事。他们就是通过这种方式来再现叙述的。他们穿上戏服，戴上面具，在舞台上用帷幔搭成一座城镇的样子，然后将那光秃秃的场景与他们逼真的动作模仿结合起来，使得观众惊为奇观——无论是扮演剧中情节的演员本身，还是帷幔，或者说他们想象中的城市。你以为我用这个比喻是要说什么？既然我们必须向来此聚会的人展示一个并非城市的东西，好像它确是一座城市，那么你愿不愿意说服自己暂时成为我们这座城市的创始人，只要在那里露个面就行？我要使这荒芜之地看起来像一座城市；现在距离对你来说并不遥远，而你将赐予的恩惠却非常巨大；因为我们希望向这里的同伴们展示自己更为光辉灿烂，如果我们除了别的装饰之外，还能以你队伍的荣光来装点自己，我们就会做到这一点。\par

\SourceRecord{Translated by H.C. Ogle.；Nicene and Post-Nicene Fathers, Second Series；Vol. 5.；Edited by Philip Schaff and Henry Wace.；Buffalo, NY: Christian Literature Publishing Co.,；1893.；<http://www.newadvent.org/fathers/291106.htm>.}

% structure: p0001 source=h1 target=section reason=page-title

\section{第七封信}

\Person{圣额我略·尼撒}\par

\Emph{致友人}\par

春天里有什么花朵如此明艳？什么鸣禽的歌声如此甜美？什么微风如此轻柔地抚平宁静的海面？什么田地如此芬芳，令农夫心旷神怡——无论它是绿意盎然的嫩苗，还是结满硕果的穗浪？同样，你书信中闪耀的光芒，如和煦的阳光，照亮了我们灵魂的春天，将我们的生命从沮丧提升为喜乐。因为，或许这样将先知的话用于我们当前的恩赐并无不妥：\Quote{在我心中众多的忧苦中，天主的安慰，}因你的仁慈，\Quote{使我的灵魂得慰藉，}如同阳光，温暖并鼓舞了我们被霜冻摧残的生命。因为两者都达到了极致——一方面，我困境的严酷；另一方面，你恩惠的甘甜。如果你仅仅通过告知我们你即将到来的喜讯，就使我们如此欢欣，以致一切都从极度的悲伤转变为光明，那么你可贵而慈祥的到来，甚至仅仅是一睹，又将成就何等的事呢？你悦耳的声音在我们耳中响起，将给我们的灵魂带来何等的安慰？愿这一切藉着天主的美善恩佑早日实现，祂赐予疲惫者痛苦中的喘息，赐予受苦者安息。但请确信，当我们审视自己的处境时，对目前的状况深感悲痛，人们不停地撕裂我们；然而，当我们转向你的卓越时，我们承认，对天主眷顾的安排怀有极大的感恩之心，因为我们能够在你的近旁享受你对我们的甜蜜和善意，尽情饱餐这样的灵性食粮，如果这类恩赐也有饱足的话。\par

\SourceRecord{Translated by H.C. Ogle.；Nicene and Post-Nicene Fathers, Second Series；Vol. 5.；Edited by Philip Schaff and Henry Wace.；Buffalo, NY: Christian Literature Publishing Co.,；1893.；<http://www.newadvent.org/fathers/291107.htm>.}

% structure: p0001 source=h1 target=section reason=page-title

\section{第八封书信}

\Emph{\Emph{圣额我略·尼撒 著}}\par

\Emph{致一位古典学学生}\par

当我依照惯例，寻找一段合适且恰当的序言——当然是指出自\WorkTitle{圣经}——置于信首时，我不知该选哪一段，并非因找不到合适的，而是因我觉得向对这事一无所知的人写这些是多此一举。因为你追求世俗文学的熱忱已无可辩驳地向我们证明，你并不关心神圣之事。因此，我將只字不提圣经经文，而要从你所喜爱的诗人那里选取一段前言，以契合你的文学品味。你们偉大教育大师笔下引入了一位老人，他在长期受苦后终于再次见到了自己的儿子，还有他的孙子，那份喜悦溢于言表。而他尤为欢欣的主题是\Person{尤利西斯}和\Person{忒勒马科斯}之间的较量，两人争夺最高的英勇奖赏；尽管他回忆自己對抗\OriginalTerm{克法利尼亚人}{Cephallenians}的功绩也增添了话语的份量。至于你和令尊，當你们像他们接待\Person{拉埃尔特斯}一样，以仁爱接待我时，你们以最为高尚的竞赛争夺美德的奖赏，尽一切可能的尊敬与善意待我们；他以多种方式——在此无需赘述——，你则以从\Place{卡帕多细亚}频频寄来的书信向我投掷。那么，我这老人当如何？我称那一日为有福的，因我亲眼目睹了这样一场父子之间的竞赛。愿你永不止歇地实现那位卓越可敬之父亲的正当祈愿，并以其继承自他的声望，在一切善工上更加踊跃。我将成为你们双方都悦纳的裁判，因我要把首奖判给你，胜过你的父亲，也把你的父亲判为胜過你。我们将栖身于粗糙的\Place{伊萨卡}——粗糙不单指石头，更指居民的习性——，那是一个有许多求婚者的岛屿，他们最觊觎的是所追求之人的财产，并借此侮辱了他们的意中人：他们以婚姻威胁她的贞洁，其行径堪似一个\Person{墨兰托}，或某个类似角色；因为那里没有一位\Person{尤利西斯}能用他的弓使他们清醒。你见我如何以老人家的方式，唠叨起与你无关的事来。但请你念及我的白发，对我多加包涵；因为爱唠叨如同老眼昏花或四肢无力一样，正是老年的特征。而你，以你敏捷活泼的言辞款待我们，恰如一位勇敢的青年，将使我们的耄耋之年返老还童，以你这般得体的悉心关怀，支撑我们漫长年岁的孱弱。\par

\SourceRecord{Translated by H.C. Ogle.；Nicene and Post-Nicene Fathers, Second Series；Vol. 5.；Edited by Philip Schaff and Henry Wace.；Buffalo, NY: Christian Literature Publishing Co.,；1893.；<http://www.newadvent.org/fathers/291108.htm>.}

% structure: p0001 source=h1 target=section reason=page-title

\section{第九封书信}

\Emph{圣额我略·尼撒}\par

\Emph{邀请}\par

春天并不习惯于一下子就展现出它灿烂的美丽，而是有春天的前奏：阳光温和地温暖着冰冻的大地表面，半掩在土块下的嫩芽，吹拂大地的微风，使空气的滋养和生殖力深深渗入。可以看到鲜嫩的草，冬天驱逐的鸟儿归来，以及许多此类征兆，这些与其说是春天本身，不如说是春天的迹象。然而这些是甜美的，因为它们是至甜之物的指示。我讲这些是什么意思呢？既然你信中传到我们这里的善意表达，作为你内在宝藏的先行者，以美好的前奏带来了我们期待从你手中得到的祝福的喜讯，我们既欢迎这些信所传达的恩赐，如同春天初绽的花朵，又祈求我们很快能享受你身上这个季节的全部美丽。因为，请确信，我们因这里人们的激情、怨恨和他们的行为而深感痛苦；就像雨水进屋后在村舍中结冰（我将从我所在地区的天气作比喻），湿气进入后，如果它覆盖在已经冻结的表面上，就会在冰上凝结，增加已有的冰块；同样，可以注意到这邻近地区大多数人的性格也有类似情况：他们总是策划和发明一些恶毒的事，新的恶行凝结在已有的恶行之上，又一个在上面，再一个，如此无休止，他们的仇恨和邪恶的增加没有止境；因此我们非常需要众多祈祷，愿\OriginalTerm{圣神}{Holy Spirit}的恩宠快快吹拂他们，消融他们仇恨的苦涩，融化他们因恶意而凝结的冰霜。因此，本来甜美的春天，对那些在经过这样的风暴之后期待你的人来说，变得比以往更加令人渴望。那么，不要让这恩赐延迟。尤其因为我们伟大的节日临近，生养你的土地为自己的宝藏而欢庆，比\Place{本都}为我们的宝藏而欢庆更为合理。那么，来吧，亲爱的，给我们带来许多祝福，甚至你本人；因为这将充满我们真福的尺度。\par

\SourceRecord{Translated by H.C. Ogle.；Nicene and Post-Nicene Fathers, Second Series；Vol. 5.；Edited by Philip Schaff and Henry Wace.；Buffalo, NY: Christian Literature Publishing Co.,；1893.；<http://www.newadvent.org/fathers/291109.htm>.}

% structure: p0001 source=h1 target=section reason=page-title

\section{\WorkTitle{书信第十}}

\Emph{尼撒的圣额我略著}\par

\Emph{致\Person{利巴尼乌斯}}\par

我曾听一位医者讲述一桩奇特的自然造化。其言如下：一人患顽疾难治，遂指责医学技艺远不如其所声称者，因凡为疗疾所设之法皆无效。其后，忽有超出其希望的佳音传来，此事竟代行医术之功，使其疾病止息。无论是因为灵魂因解脱焦虑之满溢感，并因骤然反弹，而使身体随同自身处于同样状态，抑或缘于其他途径——我无法断言，因我无暇深究此等论题，而告知此事者亦未说明缘由。但我方才恰好适时忆起此故事，窃以为甚合时宜：因当我身体欠安时——此刻我无需详述自与您分别至今所遇诸般忧烦之确切缘由——忽有人告知我，您那无与伦比的博学所写之函已到达，我一接到书信并浏览您所写之内容，立刻，首先我的灵魂受到影响，如同我在全世界面前被宣告为最辉煌成就的英雄——我如此珍视您在信中赐予我的见证——继而我的身体健康也立即开始改善：我提供了与方才所述故事相同奇事的例证：当我阅读书信前半部分时，我尚在病中，而当我阅读同一书信的后半部分时，我已痊愈。以上所述即为此事。然而，既然\Person{基内吉乌斯}是那恩惠的机缘，您因行善之丰沛富余，不仅能惠及我等，亦能惠及我等之恩人；而他正是我等之恩人，如前所述，因他成为我们得以收到您书信之原因与机缘；因此，他理应获得我们的善意。但若您问起我们的老师——倘若我们真被认为有所学习——您会发现他们是\Person{保禄}、\Person{若望}以及其余宗徒和先知；若我说得不太狂妄，声称在那您如此卓越的技艺上有所知，则公正的评判者宣称，修辞学的规则从您那里，如同从涌泉中流出，流向所有在那领域有所追求者。我曾听可敬的\Person{巴西略}向众人如此说，\Person{巴西略}，他曾是您的弟子，却是我的父亲和老师。然而，首先您要确信，我从我老师的训诲中并未获得丰裕滋养，因我只与家兄相处短暂时日，仅从他神圣的舌端得到刚好足以分辨那些未受修辞学启迪之人的无知；其次，每当我得闲时，我便将时间与精力投入此学，因而迷恋您的美丽，却从未获得我热爱的对象。若然，若我们从未有过一位老师（我认为我们确是如此），且又若将您对我们的看法视为不真乃是失当——不，您的陈述是正确的，在您判断中我们并非全然卑微——请允许我斗胆将我们可能达到的任何造诣之原因归于您。因为若\Person{巴西略}是我们修辞学的作者，且他的财富来自您的宝库，则我们所拥有的乃属于您，即使我们通过他人接受。但若我们的成就甚少，好比瓮中之水，但它仍来自\Place{尼罗河}。\par

\SourceRecord{Translated by H.C. Ogle.；Nicene and Post-Nicene Fathers, Second Series；Vol. 5.；Edited by Philip Schaff and Henry Wace.；Buffalo, NY: Christian Literature Publishing Co.,；1893.；<http://www.newadvent.org/fathers/291110.htm>.}

% structure: p0001 source=h1 target=section reason=page-title

\section{第十一封信}

\Emph{\Person{圣尼撒的额我略} 著}\par

\Emph{致\Person{利巴尼乌斯}}\par

\Person{罗马人}依祖传习俗，在冬季白日渐长、太阳升至天空高处之时，举行节庆。月初被视为神圣，他们藉此日占卜整年之吉凶，致力于预测幸运、喜乐与财富。我为何如此开篇？因为我亦庆祝此节，且与众人一样得到了金礼。那时，也正如他人一般，有黄金落入我手中——并非凡俗之金，即权贵所珍藏、拥有者所馈赠的沉重、卑贱且无灵之物；而是超越一切财富的，如\Person{品达}所言，在明理之人眼中最为珍贵的礼物——即您的信札及其所包含的无尽财富。事情是这样的：那天，我在前往\Place{加帕多家}首府途中，遇见一位熟人，他将您的信作为新年礼物交给了我。我欣喜若狂，将这份宝藏向所有在场之人敞开；大家都分享了它，每人尽得全部，毫无争执，而我亦未减少分毫。因这信札犹如宴席的入场券，经众人之手传递后，成为每个人的私有财富——有人凭持续诵读将字句铭刻于记忆中，有人则将其摹写于板片之上——而它又重新回到我手中，带给我比贵金属予富人眼中的喜悦更大的欢乐。既然，即使对农夫——打个粗浅的比方——对他们已完成劳动的赞许，也是后续劳动的强劲动力，那么，请容我们以您所给予的作为种子，写信以激发您的回信。但我向您恳求一项公众普世之恩惠：请您不再保留您在信末含糊暗示的意图。因为我认为，因有些人背弃希腊语而投奔蛮族，成为雇佣兵并选择兵士口粮而非辩才之荣誉，便因此全盘否定修辞学，并判人类生活如野兽般无声，这绝非公平的决定。若您对修辞学实施此严厉判决，谁还敢开口呢？但也许最好提醒您圣经中的一段话。因我们的\OriginalTerm{圣言}{Word}命令那些有能力行善者，不要看受惠者的态度，以致只热心施恩于知恩者，而对忘恩者关闭善门；却要效法万有的安排者，祂将受造的美物平等地分施给所有人，无论善恶。鉴于此，可敬的先生，请您在生活中展现如过去岁月所显示的那样的您。因为看不见太阳的人并不因此阻碍太阳的存在。同样，也不应因那些灵魂感知上短视的人而暗淡您口才的光芒。至于\Person{基内吉乌斯}，我祈祷他尽可能远离如今困扰年轻人的通病，并自愿投身修辞学的研究。但若他另有所图，即使他不情愿，也应当强迫他去学习，以避免像先前放弃修辞学追求者现在所处的不幸与可耻的境地。\par

\SourceRecord{Translated by H.C. Ogle.；Nicene and Post-Nicene Fathers, Second Series；Vol. 5.；Edited by Philip Schaff and Henry Wace.；Buffalo, NY: Christian Literature Publishing Co.,；1893.；<http://www.newadvent.org/fathers/291111.htm>.}

% structure: p0001 source=h1 target=section reason=page-title

\section{\WorkTitle{Letter 12}}

\Person{圣国瑞·尼撒}\par

\Emph{论其驳斥欧诺弥著作}\par

我们卡帕多细亚人几乎在一切令拥有者幸福的事上都贫乏，但最匮乏的莫过于能书写的人。请相信，这就是我迟迟未给你寄信的原因：虽然我对（欧诺弥）异端的答辩早已完成，却无人抄写。正是缺乏缮写者，才令我们蒙受迟延或无力作答的嫌疑。但如今，感谢天主，缮写者和校对者总算来了，我便将此论著寄给你；并非像依索克拉底所说，作为赠礼——因我认为它尚不值得以实礼相待——而是为了激发青年人的热情，凭借鼓舞早年生来的乐观气质，使他们能奋勇与仇敌作战。若书中某部分值得深思，经检验某些段落（尤其是那些置于\Quote{试炼}之前的序言及类似文风，或许也包括书中某些教义部分），你会觉得其撰写并非毫无可取。但无论你得出何种结论，你自会将其读给那些不像玩球者的人，如同读给导师和校正者：玩球者分站三处，互相抛接，准确瞄准，接住来自各方的球；他们阻挠中间球员跳跃接球，佯装要以某种表情和手势向左右投掷，见其向某方疾冲，便反方向掷球，以诡计欺骗其预期。我们多数人如今正是如此：放下严肃目的，以宽容者为戏，如同玩球般对待他人，而不兑现所许诺的有利希望，以险恶的结局诱骗信任我们者们的灵魂。和解的信函、爱抚、信物、赠礼、书信中亲切的拥抱——这些都像是佯装向右掷球。但人所期待的快乐却变成指控、阴谋、诽谤、诋毁、控诉、断章取义的句子被拾起，转而伤害人。你在希望中是有福的，因你历经这一切试炼而信赖天主。但我们恳求你，不要看我们的话语，而要看主在福音中的教导。对痛苦中的人，还有谁比他更痛苦？他的痛苦能帮助不幸之人获得适当结局吗？\BibleQuote{正如祂说：上主要报复，我要报复，上主说。}但你，最好的人啊，请继续以与你相称的方式前行，信赖天主，不要因我们不幸的景象而阻碍你行善持真，而要将合适而公正的结局交托给那按正义审判的天主，并随从神圣智慧引导你行事。的确，约瑟最终没有因兄弟们的嫉妒而悲伤，因为亲族的恶意反而成了他通往王权的道路。\par

\SourceRecord{Translated by H.C. Ogle.；Nicene and Post-Nicene Fathers, Second Series；Vol. 5.；Edited by Philip Schaff and Henry Wace.；Buffalo, NY: Christian Literature Publishing Co.,；1893.；<http://www.newadvent.org/fathers/291112.htm>.}

% structure: p0001 source=h1 target=section reason=page-title

\section{第十三封书信}

\Emph{\Person{圣额我略·尼撒} 著}\par

\Emph{致\Place{尼苛默底亚}教会}\par

愿仁慈的圣父、一切安慰的天主，祂以智慧为最佳安排万有，藉祂的恩宠眷顾你们，并亲自安慰你们，在你们内成就祂所喜悦的事。愿我们的主耶稣基督的恩宠、圣神的共融临于你们，使你们获得一切忧苦与患难的医治，并在一切善事上进步，为教会的成全、你们灵魂的建树，并赞美祂光荣的圣名。但在此，我们愿在你们的爱德前为自己辩护：我们并未疏忽对所托付之职责的交待，无论是过去，还是自从可纪念的\Person{帕特里西乌}离世之后。但我们坚持我们的教会有许多困难，我们身体的力量也大大衰退，随着年岁增长，这是自然的；而你们的阁下对我们也很怠慢，因为没有一封来信劝勉我们承担此任务，你们教会与我们之间也未曾保持任何联系，尽管你们可纪念的主教\Person{欧弗拉西乌}曾以极大的圣洁，用爱如锁链般将我们的卑微与他及你们联结在一起。但即使以前，或因我们未照管你们，或因你们的虔诚未鼓励我们，这爱的债务未能清偿，但现在无论如何，我们向天主祈祷，将你们的祈祷作为盟友与我们的愿望结合，愿我们能尽快访问你们，与你们一同受安慰，并与你们一同殷勤行事，如主所引导我们，以便找出纠正已发生的混乱的方法，并确保未来的安全，使你们不再因这不和而分散，一人向此方向脱离教会，另一人向彼方向脱离，从而成为魔鬼的笑柄，魔鬼的愿望和事业（与天主的旨意直接相反）是无人得救，也无人认识真理。弟兄们，你们想，当我们从那些向我们报告你们情况的人那里听到，你们没有回归更好的事，而是那些一度偏离之人的决意始终沿同一进程；并且——如同水从水渠常溢出邻近的堤岸，斜向流走，除非堵住漏洞，否则几乎无法将其引回水道，因为被淹没的地面已随水流冲成沟壑——同样，那些离开教会之人的进程，一旦因个人动机从正直和正确的信仰偏斜，就已深陷习惯的轨迹，不易回归曾有的恩宠。因此，你们的事务需要一位明智而坚强的管理者，他善于正确引导这些任性的性情，以便能够将这水流无序的循环召回其原有的美丽，使你们虔诚的田地因和平的灌溉之流再次丰茂。为此原因，你们全体迫切需要极大的热忱和炽热的愿望，以便由圣神指定这样一位作你们的首领，他独眼注视天主的事，不向左或向右看任何人所追求的事物。因为，我想，古代的法律没有给肋未人在土地总继承中的份额，正是为使他唯独以天主为他的产业，常专注于自身的产业，而不着眼任何物质对象。\par

\EditorialNote{以下内容无法理解，可能有所缺失。}\par

因为庶民不应干预与他们无关、而只属于他人的事务。因为你们每人应各管各的事，这样最有益的事才能发生[好使你们的教会再度繁荣]，当那些分散的人重新回到一个身体的合一，并藉着虔诚光荣天主的人建立起属灵的平安。为此，我认为，在你们的选举中寻找高尚的资格是好的，使被任命为主席的人适合该职位。现在，宗徒的训令并不教导我们在主教的德性中看重出身高贵、财富和世俗眼光的显赫；但如果这一切未经追求而伴随你们的属灵领袖，我们也不拒绝，不过只视之为偶然跟随身体的影子；我们仍会欢迎更珍贵的禀赋，即使它们恰好没有那些命运的恩赐。\Person{亚毛斯}先知是个牧羊人；\Person{伯多禄}是渔夫，他的兄弟\Person{安德肋}也从事同样的职业；崇高的\Person{若望}也是如此；\Person{保禄}是制帐篷的，\Person{玛窦}是税吏，其余宗徒也一样——不是执政官、将军、总督，也不是以修辞和哲学闻名，而是贫穷的，没有学问，却从生活中更卑微的职业开始：然而他们的声音却传遍大地，他们的话语传到地极。\BibleQuote{弟兄们！你们要审视自己的蒙召：按人的标准，蒙召的有智慧的并不多，有能力的也不多，有显贵身份的也不多；但天主却拣选了世上愚拙的\BibleRef{格前 1:26}。} 也许即使在现在，如人们所见，当一个人因贫穷而不能多做，或因出身寒微（而非品性）而受到轻视时，这也被视为愚拙。但谁知道恩宠的傅油之角是否倾注在那人身上，即使他不如崇高显赫的人？什么更有利于\Place{罗马}教会：在开始的时候由某个出身高贵、傲慢的元老主持，还是由渔夫\Person{伯多禄}、这位没有任何世俗优势吸引人的人主持？他有什么房子、什么奴隶、什么财产以源源不断的财富提供奢华？但那个陌生人，没有桌子，没有屋顶，却比拥有一切的人更富有，因为一无所有而完全拥有天主。同样，\Place{美索不达米亚}的人民，尽管他们中间有富有的总督，却更喜欢\Person{多默}超过所有其他人来主持他们的教会；\Place{克里特}人喜欢\Person{弟铎}，\Place{耶路撒冷}的居民喜欢\Person{雅各伯}，而我们\Place{卡帕多细亚}人喜欢那位在十字架前承认主的天主性的百夫长，尽管当时有许多家世显赫的人，他们的财富足以维持马厩，并以在元老院中占有首位而自豪。在所有教会中，我们可以看见那些按天主标准为大的胜过世俗的显赫。我想，你们在现今也应当着眼于这些属灵的资格，如果你们真有意复兴教会的古代荣耀。因为没有什么比你们自己的历史更清楚：古时，在你们附近的城市繁荣之前，政府所在地曾是在你们那里，在\Place{彼提尼雅}的城市中没有比你们的城市更优越的。而现在，诚然，曾经装饰它的公共建筑已经消失，但那由人构成的城市——无论我们看数量还是质量——正迅速上升到昔日的辉煌水平。因此，你们应当怀有与你们现今所拥有的祝福高度相称的思想，并在你们面前的工事中激发与你们城市壮丽相应的热情，好使你们能找到一位管理平信徒的人，证明自己不辜负你们。因为，弟兄们，这是可耻且极其怪异的：既然没有人会不精通航海就当舵手，那坐在教会船舵上的人却不知道如何将与他同航的灵魂安全地带入天主的港口！多少次教会连同人员的沉船事故因领袖的无经验而发生！谁能计算有多少灾难本可避免，如果那些掌舵者有任何舵手的技能？不，我们委托铁匠打造器皿，不是给那些不懂此事的人，而是给熟悉铁匠艺术的人；那么我们不更应该把灵魂托付给那位善于以圣神的热火软化它们、并藉理性工具的印记将你们每个人塑造为合用的器皿的人吗？为此，受默感的宗徒在\BibleBook{弟茂德}{前书}中吩咐我们思考\BibleRef{弟前 3:2}，对所有听到的人发出训令，他说主教必须无可指责。宗徒所关心的就只是这些吗，让被提升到司铎职位的人无可指责？还有什么比把一切可能的资格都集中在一人身上更大的好处？但他很清楚，属下会被上级的性格所塑造，向导的正直行走也会成为随从者的正直行走。因为师傅如何，他也使门徒成为如何。因为不可能学过铁匠手艺的人去从事纺织，或者只学过织布的人变成演说家或数学家：相反，门徒在师傅身上看到的，他采纳并转移到自己身上。为此，圣经说：\BibleQuote{凡完全的门徒，都要如同他的师傅。} 那么，弟兄们，一个人有可能品性谦卑温和、有节制、超脱贪财、智慧于神圣之事、在言行上受过美德和体谅的训练，而没有在师傅身上看到这些品质吗？不，我不知道一个人如何在世俗的学校做门徒而能成为属灵的人。因为那些努力效法师傅的人岂能不肖似他？对于口渴的人，水渠的壮丽有何益处，如果里面没有水，即使各种形状的柱子对称排列高高撑起门楣？口渴的人更愿意选择什么来满足自己的需要：是看到优美排列的大理石，还是找到好的泉水，即使它流经木管，只要它流出的溪水清澈可饮？同样，弟兄们，那些追求敬虔的人应当忽视外在装饰的炫耀，无论一个人是得意于有权势的朋友，还是以长长的头衔名单自夸，或是夸耀每年收到大量收入，或是因高贵的祖先而自傲，或是心意四围被自负的烟雾笼罩，都不应与这样的人有任何关系，就像与干涸的水渠无关一样，如果他在生活中没有显示出高位所需的根本而不可或缺的品质。但是，用圣神的灯进行搜寻，你们应尽可能寻求\BibleQuote{一座关闭的花园，一个密封的泉源\BibleRef{歌 4:12}}，好使藉由你们的选举，那喜乐的花园被打开，泉源的水被释放，成为天主教会共同的财富。愿天主赐予在你们中间很快找到这样一个人，他将是合用的器皿，教会的柱石。但我们信赖主，事情会如此，如果你们藉由同心合意的恩宠意愿看见那善的，以主的旨意胜过你们自己的意愿，以及那在祂眼中被认可、完全和悦乐的事，好使在你们中间有如此蒙福的结局，使我们可以欢乐，你们可以夸胜，万有的天主、光荣永远属于祂的，得受光荣。\par

\SourceRecord{Translated by H.C. Ogle.；Nicene and Post-Nicene Fathers, Second Series；Vol. 5.；Edited by Philip Schaff and Henry Wace.；Buffalo, NY: Christian Literature Publishing Co.,；1893.；<http://www.newadvent.org/fathers/291113.htm>.}

% structure: p0001 source=h1 target=section reason=page-title

\section{第十四封信}

\Emph{尼撒的圣国瑞著}\par

\Emph{致梅利特内的主教}\par

美丽对象的肖像，当它们清晰保存原美的印记时，何其美丽！因为在你那真正美丽的灵魂上，我从你信中的甘甜看到了极其清晰的形象：正如福音所说，你以蜜充满了\Quote{出自心的丰盈}。为此，我仿佛亲眼见到了你，并从你信中流露的情感中享受了你令人振奋的陪伴；我时常将你的信拿在手中，从头到尾反复阅读，却只是更加强烈地渴望那享受，毫无餍足之感。这种情感既不能终结我的喜乐，正如任何本质上美丽珍贵的事物带来的喜乐也无法终结一样。因为我们既因恒常参与恩惠而未曾削弱对瞻仰太阳的渴望，也因不间断地享有健康而未曾停止期望其延续；同样，我们深信，对你美善的享受——我们曾多次面对面体验，如今又通过书信体验——永远不会达到餍足。但我们的情况如同那些因某种情况而患上无法止息之渴的人；因为同样，我们越品尝你的仁慈，就越口渴。但除非你以为我们的话只是甜言蜜语和不实的奉承——你肯定不会这样想，因你在各方面都是如此，尤其对我们，若有人曾如是，你更是善良而坚定——你必定会相信我的话：你信中的恩宠，如同某种药方敷在我的眼上，止住了我永流的\Quote{泪泉}；我们将希望寄托于你圣善祈祷的良药，期待不久之后，我们灵魂的疾病将完全痊愈：尽管目前，我们处于这样的状况，以致我们不愿让爱我们的人耳闻，而将真相埋于沉默，以免使那些忠诚爱我们的人分担我们的烦恼。因为当我们想到，我们失去了最亲爱的人，陷于战争之中，而且我们的孩子——那些我们曾蒙恩在灵性痛楚中为天主所生、因爱之律法与我们紧密相连的孩子，在他们自己受考验的苦难时刻也曾向我们表达关爱——我们被迫将他们留在身后；此外，还有亲爱的家园、兄弟、亲属、同伴、密友、朋友、炉灶、餐桌、地窖、床榻、座位、行囊、交谈、眼泪——这些是何等甘甜，因长久习惯而何等珍贵，我无需向深知这一切的你多写——但为了不使你更厌烦，你自己想想我用这些福分换来了什么吧。如今我生命将尽，却开始重新生活，被迫学习如今流行的灵活多变之性格：但在狡猾之流变的学校里，我们乃是晚学；因此我们常因自己在学习这门新学问时的笨拙和不适应而脸红。但我们的对手，装备了这智慧的一切训练，既能守住已学的，又能发明未学的。他们的作战方法便是远处骚扰，然后按约定信号组成密集方阵；他们先以对自己有利的话作为开场，用夸张手法发动突袭，从四面八方召集盟友。然而，一股巨大而不可战胜的狡诈伴随着他们，先行于前领导他们的队伍，如同左右手俱能作战的斗士，在军队前方双手搏击：一边向臣民征收贡品，一边击打挡路者。但若你愿意询问我们内部事务的状况，你会发现其他相称的烦恼：一间令人窒息的茅屋，充满寒冷、阴暗、局促和一切此类便利；一种成为众人挑剔观察目标的生活：声音、目光、披风穿戴方式、手的动作、脚的位置，以及一切，都成为好事者的谈资。除非一个人不时发出深沉的呼吸，除非呼吸中伴随着持续的呻吟，除非上衣优雅地穿过腰带（更不必说腰带本身的不常用），除非我们的披风斜挂背后——这些细节中任何一项的疏忽，都成为向我们开战的借口。基于这样的理由，他们聚集起来攻击我们，逐人、逐镇，甚至至各种偏远之处。唉，人不可能总是顺利或总是倒霉，因为每个人的生活都由相反之事组成。但如果赖天主的恩宠，你的帮助能坚定地支持我们，我们将承受这许多烦恼，以期望始终分享你的美善。因此，愿你永不停止赐予我们这样的恩惠，藉以振奋我们，并为你自己更丰厚地预备遵守诫命者所应许的赏报。\par

\SourceRecord{Translated by H.C. Ogle.；Nicene and Post-Nicene Fathers, Second Series；Vol. 5.；Edited by Philip Schaff and Henry Wace.；Buffalo, NY: Christian Literature Publishing Co.,；1893.；<http://www.newadvent.org/fathers/291114.htm>.}

% structure: p0001 source=h1 target=section reason=page-title

\section{\WorkTitle{第十五封信}}

\Signature{由\Person{圣额我略·尼撒}著}\par

\Emph{致律师\Person{阿德尔斐乌斯}}\par

我从神圣的\Place{瓦诺塔}写这封信给你，如果我用它的本地称呼而亏待了这地方，那便是我的不对——我说亏待，是因为就其名称而言，它并不优雅。同时，这地方的美景虽大，却非这个\Place{加拉提亚}的称号所能传达：需要亲眼目睹才能领略其美。因为，虽然我此前曾在许多地方见过许多事物，也通过古代作家的叙述以言辞描述观察过许多事物，但与我在此所见的可爱相比，我既认为我所见的一切，也认为我所闻的一切，都不足道。你的\Place{赫利孔山}算不得什么：\Place{幸福岛}只是传说：\Place{西库昂平原}微不足道：关于\Place{佩纽斯河}的叙述又是诗意的夸张——他们说那条河以其丰沛的水流泛滥过两岸，为\Place{色撒利人}造就了著名的\Place{坦佩谷}。哎呀，在我所提及的这些地方中，有哪一个能像\Place{瓦诺塔}那样向我们展示其自身的美呢？因为若有人在此寻求自然之美，它无需任何人工装饰：若有人考虑人工辅助对它所做的，则所做甚多且甚佳，足以克服甚至自然的不利条件。大自然以质朴的优美装点大地，赐予此地的恩赐如下：下方，\Place{哈吕斯河}以其河岸使这地方悦目，并像一条金色丝带在深紫色中闪闪发光，河水因冲刷土壤而泛红。上方，一座密布树木的山脉绵延着长长的山脊，处处覆盖着橡树的枝叶，值得找一位\Person{荷马}来歌颂它，胜过那位诗人称为\Quote{远见摇曳叶}的\Place{伊萨卡}的\Place{内里图斯山}。但天然生长的林木沿山坡而下，在山脚与人类的耕作种植相遇。因为立刻，葡萄藤遍布山坡、隆丘和山脚的低洼处，用它们的色彩像一件绿色斗篷覆盖所有较低的地面；此时的季节甚至更添其美，展现出令人惊叹的葡萄串。确实，这更令我惊奇，当邻近地区果实尚未成熟时，在此处却可享用饱满的葡萄串，饱享它们的完美。然后，远处，像某座大灯塔的守望之火，建筑的美丽在我们眼前闪耀。在我们进入时，左侧是为殉道者建造的小教堂，结构尚未完成，仍缺屋顶，但尽管如此，已显得颇为可观。正前方路上是房屋的美丽，各部分以某种精巧的设计彼此区分。有突出的塔楼，以及在高大宽阔的拱形树列间为宴会所做的准备，这些树列冠于门前的入口。然后建筑四周是\Place{费阿刻斯人}的花园；不如说，不要让\Place{瓦诺塔}的美景与它们比较而受辱。\Person{荷马}从未见过\Quote{果实鲜亮的苹果}像我们这里的这般，其色泽之鲜艳近乎自身花朵的色调：他从未见过梨子比新抛光的象牙更白。对于桃子的各种品种，多样而多形，却又由不同种类混合而成，又能说什么呢？正如那些描绘\Quote{羊鹿}和\Quote{半人马}等形象的人，将不同种类混合，自认为比自然更聪明，这种水果也是如此：在艺术专制下，自然将一个变成杏仁，另一个变成核桃，又一个变成\OriginalTerm{多拉基努斯}{Doracinus}，名称和风味都混合了。在所有这些中，单棵树木的数量比其美更引人注目；然而它们展现出种植的雅致布局，以及那种和谐的绘图形式——我称之为绘图，因为这种奇迹更属于画家的技艺而非园丁。自然如此乐于迎合那些安排这些装置者的设计，以至于似乎无法用言语表达。谁能找到合适的词来描述攀爬的葡萄藤下的道路、葡萄串的甜美荫凉、以及那新奇的墙结构——玫瑰的嫩枝和葡萄藤的卷须交织在一起，形成一道抵御侧面攻击的防御工事，还有这条小路顶端的水池和在那里养殖的鱼呢？关于所有这些，负责大人府邸的人们以某种天真的和善乐意担任我们的向导，并向我们指出它们，向我们展示你费心打理的每一件东西，仿佛通过我们，他们是在向你本人表示礼貌。在那里，一个少年像魔术师一样，向我们展示了一个在自然中并不常见的奇观：他潜入深水，随意捞起他选中的鱼；这些鱼似乎并不陌生于渔夫的触碰，在艺术家手下温顺服从，像训练有素的狗。然后他们带我到一个房子状的地方休息——我称之为房子，因为入口显示如此，但当我们进入时，迎接我们的不是房子，而是一个柱廊。柱廊高高矗立在一个深水池上方：支撑柱廊的底座呈三角形，像一道通向内部乐趣的大门，被水冲刷着。正前方内部，一种房屋占据了三角形的顶点，屋顶高耸，四周阳光照耀，装饰着各种绘画；以至于这个地方几乎让我们忘记了先前的所见。房屋吸引着我们；而水池上的柱廊又是一道独特的风景。因为出色的鱼会从深处游到水面，像有翼之物跃入空中，仿佛故意嘲笑我们这些旱地生物。它们展示半身，在空中翻滚，然后再次潜入深处。其他鱼群则成队列依次跟随，对于不习惯的眼睛是一种景象：而在另一处，可以看到另一群鱼簇拥在一块面包屑周围，彼此推挤，这里一条跃起，那里一条下潜。但即便是这些，也被送来的葡萄——装在用枝条编成的篮子里——丰盛的时令水果、早餐的准备、各种美味佳肴、甜品、祝酒和酒杯所使我们忘却。所以，既然我已饱足且昏昏欲睡，我让一位抄写员坐在我身旁，像做梦一样将这篇喋喋不休的信寄给你的口才。但我希望向你本人和你的朋友们完整地叙述你家的美景，不是用纸和墨，而是用我自己的声音和舌头。\par

\SourceRecord{Translated by H.A. Wilson.；Nicene and Post-Nicene Fathers, Second Series；Vol. 5.；Edited by Philip Schaff and Henry Wace.；Buffalo, NY: Christian Literature Publishing Co.,；1893.；<http://www.newadvent.org/fathers/291115.htm>.}

% structure: p0001 source=h1 target=section reason=page-title

\section{第十六封书信}

\Emph{\Person{圣尼撒的额我略} 著}\par

\Emph{致\Person{安斐洛希}}\par

我深信，藉天主的恩宠，殉道者圣堂的事务进展顺利。惟愿您对此事有意。我们手头的工作将因天主的德能而完成，祂无论在哪里发言，都能将言语化为行动。既然，如宗徒所说，\BibleQuote{那在你们内开始了美好工作的，必予以完成}，我劝您在此事上也成为伟大的保禄的效法者，将我们的希望推进至实际完成，并为我们送来足够多的工作者，以完成我们手头的工作。\par

阁下或许可通过计算整体工程将达到的尺寸而得知：为此，我将尽力用语言描述整个结构。圣堂的形状是十字形，其图形由四个结构完整构成，正如您所料。建筑物的连接处相互交错，如同我们在十字形图案中常见的那样。但在十字形内部有一个圆，由八个角分隔（我称八角形图形为圆，是因它的圆周），这样八角形相对的两对边通过拱将中央圆与相邻的建筑块连接起来；而八角形其余四条边，位于四边形建筑之间，它们本身不会延伸到建筑处，而是每条边上将描绘一个半圆，状如贝壳，顶部以拱终止：这样总共有八个拱，通过它们，四边形和半圆形建筑侧向地与中央结构连接。在由角形成的砌体中，将有相等数量的柱子，既为装饰也为坚固，这些柱子将再承载与内部同等大小的拱。在以上八个拱之上，对称地加上一层窗户，八角形建筑将升至四肘高；其上部分将是一个圆锥形，状如陀螺，因为拱顶将屋顶图形从全宽收窄至尖楔。下面的尺寸将是：每个四边形建筑的宽度为八肘，长度为其一倍半，高度等于宽度比例所允许的。半圆形建筑也是如此。墩柱之间的全长同样延伸至八肘，深度将由圆规固定于边的中点并向端点延伸所画的弧线决定。此情形下，高度也将由与宽度的比例决定。墙壁的厚度，从内部测量的这些空间向内间隔三英尺，将环绕整个建筑。\par

我以这严肃的琐事烦扰尊座，是为此目的：藉由墙壁的厚度和中间的空隙，您可以精确地推算出若干尺数所构成的总量；因为您的才智在所有事务中都极其敏锐，且因天主的恩宠，在您所愿的任何主题上都能畅行无阻，因此您有可能通过精妙的计算，确定所有部分的总和，以便派给我们的石匠不多不少，恰合所需。我恳请您特别留意这一点：其中有些人须技艺娴熟，能建造无支撑拱顶；因为我听说如此建造比依靠支柱的更耐久。正是木材的匮乏迫使我们采用以石料覆盖整个建筑的方案；因为此地不产用于屋顶的木料。请您确信无误，因为此地有些人向我承包，以一块\OriginalTerm{斯塔特}{stater}的价格提供三十名工人，当然是指加工过的石料，且每块斯塔特附带定额口粮。但我们的石料并非此类，而是用黏土和随意石块制成的砖，因此他们无需费时精确拼接石面。我知道，就技艺和工价的公平而言，您那边的工人比本地的更符合我们的需求。雕刻工作不仅在于八根立柱本身，这些立柱需要改善和美化，而且还需要祭台式基座饰面，以及按科林斯式雕刻的柱头。门廊也将用适当装饰的大理石制成。安装其上的门扇将饰有通常用于檐口突出处进行美化的一些图案。当然，这些材料的全部由我们提供；艺术将赋予材料以形式。此外，柱廊中将有不下四十根立柱：这些也将是雕刻的石料。现在，如果我的说明已详细解释了该工程，我希望能使您的圣座在了解所需后，完全解除我们在工人方面的忧虑。然而，如果工人倾向于与我们达成有利的协议，那么如果可能，请为每日规定明确的工作量，以免他无所事事地消磨时间，然后虽无工作成果，却以为我们工作了这么多天为由要求支付工资。我知道，在多数人看来，我们在合同上如此斤斤计较，就像个讨价还价者。但恳请您原谅我；因为我常严厉斥责的那个\OriginalTerm{玛门}{Mammon}，终于尽可能远地离开了我，我想，是因我不断对他所说的胡言乱语而厌恶，并藉一道不可逾越的鸿沟——即贫穷——巩固了自己，以至于他不能到我这里来，我也不能到他那里去。这就是我坚持工人公平的原因，以便我们能完成面前的任务，不被贫穷——那可赞美且可欲的恶——所阻碍。好吧，这一切都掺杂了某种玩笑。但您，天主的人，请以可能且合法的方式，在与工人谈判时大胆承诺，他们都将从我方得到公平对待和全额工资：因为我们必会全部给付，分毫不留，正如天主也藉由您的祈祷，为我们敞开祂祝福之手。\par

\SourceRecord{Translated by H.A. Wilson.；Nicene and Post-Nicene Fathers, Second Series；Vol. 5.；Edited by Philip Schaff and Henry Wace.；Buffalo, NY: Christian Literature Publishing Co.,；1893.；<http://www.newadvent.org/fathers/291116.htm>.}

% structure: p0001 source=h1 target=section reason=page-title

\section{第十七封书信}

\Emph{\Person{圣额我略·尼撒}作}\par

\Emph{致欧斯塔西亚、安布罗西亚与巴西利莎。}\par

致最明智而虔诚的姊妹欧斯塔西亚与安布罗西亚，以及最明智而高贵的女儿巴西利莎：\Person{额我略}在主内问候。\par

与良善及蒙爱者的会面，以及主对我们人类那无限之爱的纪念标志——这些在你们所在之处显明——成为我最深切喜乐与欢欣的泉源。的确，这些事在神圣的节庆日子里加倍照耀：既因目睹那赐予我们生命之天主的救赎标记，也因与那些灵魂相遇——在这些灵魂中，主恩宠的标记以如此清晰的方式属灵地被辨明，使人可以相信\Place{白冷}、\Place{哥耳哥达}、\Place{橄榄山}与复活之地确实存在于那含蕴天主的心中。因为当基督通过良好的良心在某人内形成时，当某人因敬畏天主而钉死了肉身的冲动，并与基督同钉在十字架上时，当某人从自己身上滚开这世界幻象的重石，并从身体的坟墓中走出，开始好似在生命的新颖中行走，抛弃这人类生命的低洼山谷，以高飞的渴望攀登那充满崇高思想的天乡，那里有基督，不再感到身体的负担，反而通过贞洁将其提升，以致肉身如云般轻盈地伴随着上升的灵魂——这样的人，在我看来，应被列在那些著名人物之列，在他们身上可见主对我们人类之爱的纪念标志。因此，当我不仅以视觉感官目睹了那些圣地，而且也在你们自己身上看到了与它们相似的标志时，我被如此巨大的喜乐充满，以致这福分的描述超乎言语之能。然而，因为人类享有不掺杂任何邪恶的福分是困难的，更不用说不可能了，所以在我所尝的甘甜中混入了一些苦涩：在享受了那些福分之后，我在返回故乡的旅途中感到忧伤，此时我体会到主话语的真实：\Quote{整个世界都卧在邪恶者手下}（\BibleRef{若一 5:19}），以致可居住大地的任何部分都没有分毫不道德。因为如果那曾接受生命本身足迹的地方本身也没有清除邪恶的荆棘，我们当如何看待其他地方呢？——在那里，与福分的共融仅通过听闻和宣讲而被灌输。我为何说这话，无需言辞进一步解释；事实本身比任何言语（无论多么清晰）都更响亮地宣告这可悲的真理。\par

我们生命的立法者只命令我们有一种憎恨，我指的是对那蛇的憎恨：祂吩咐我们运用这种憎恨的能力，除了作为对抗邪恶的武器，别无其他目的。祂说：\Quote{我要把敌意放在你和它之间。}由于邪恶是复杂而多样的，圣言用蛇来寓意它，蛇那密集的鳞片象征着这种邪恶的多形性。而我们通过执行我们仇敌的旨意与这蛇结盟，从而将这憎恨转向彼此，也许不仅转向我们自己，也转向那颁布诫命者；因为祂说：\Quote{你应爱你的近人，恨你的仇敌}，命令我们只以人类之敌为仇敌，并宣告所有分享人性的人都是我们每个人的近人。但这心肠粗鄙的时代使我们与近人分离，使我们欢迎那蛇，并陶醉于它那斑驳的鳞片。因此，我断言，憎恨天主的仇敌是合法的，这种憎恨是悦乐我们主的：我所说的天主的仇敌是指那些否认我们主之荣耀的人，无论是犹太人，还是彻底的偶像崇拜者，或是那些通过\Person{亚略}的教导将受造物神化，从而采纳了犹太人的错误的人。如今，当圣父、圣子及圣神以正统的虔敬被那些相信在分明而不混乱的圣三中有一个本体、荣耀、王权、能力与普遍宰治的人们所光荣与朝拜时，在这种情况下，还有什么合理的争战借口呢？当然，当异端观点占上风时，与权威——仇敌的事业藉之得到加强——交锋是好的；那时是恐怕我们得救的教义被人类统治者压倒。但现在，当正统信仰从天的这一端到另一端被传遍全世界时，那与宣讲它的人争战的人，并非与他们争战，而是与那被宣讲者争战。确实，一个对天主有热忱的人，其目标除了以一切可能的方式宣告天主的荣耀之外，还能是什么呢？因此，只要那独生子以全心、全灵、全意被朝拜，被相信在一切方面都与圣父相同，同样，圣神也以同等的朝拜被光荣，那么这些过分精明的争辩者——他们撕裂那无缝的长袍，将主的名分给\Person{保禄}和\Person{刻法}，毫不掩饰地厌恶与敬拜基督的人接触，几乎等于在说：\Quote{离开我，我是圣的！}——还留下什么合理的争战借口呢？\par

即使他们相信自己获得的知识比他人稍多，但他们所能拥有的，岂不也只是相信\Quote{真天主之子是真天主}吗？因为在那\Quote{真天主}的信条中，已包含了所有正统思想，所有关乎我们救恩的思想。它包含了祂的良善、公义、全能：祂不容变化或改易，而是恒常不变；不能变得更坏或更好，因为前者非其本性，后者祂亦不接纳；因为有什么能高于至高者，有什么能优于至善者？事实上，祂如此与一切完美相连，且就一切改变而言，是不可改变的；祂并非偶然才显示这一属性，而是始终如此——无论是在那使祂成为人的安排之前、其间或之后；在祂为我们所做的一切行动中，祂从未将那不变不改本性的任何部分降低到与之不相称的地步。本质上不朽不变的东西总是如此；当它因安排而进入低等事物时，并不随其变化而变化；例如太阳，当它把光束投入黑暗时，并不会使那光束变暗；相反，黑暗因那光束变为光明；同样，真光照耀在我们的黑暗里，它本身并未被那黑暗所遮蔽，反而藉着自身照亮了黑暗。那么，因我们的人性处于黑暗中，正如经上所载：\Quote{他们不知道，也不明白，仍在黑暗中行走}，这黑暗世界的照明者便将祂天主性的光芒射入我们全部本性的复合体中——即透过灵魂，也透过身体——以此藉着自己的光芒完全占有人性，并将它提升，使它成为祂自己所是的那一位。正如这天主性虽居于可朽的身体，却并未变得可朽；同样，它虽医治我们灵魂中易变的部分，却并未朝任何变化的方向改变：在医学中，身体的医师在接触病人时，非但自己不会染上疾病，反而藉此完全治愈了患部。同样，谁也不可误解福音的话，以为基督内的人性是藉着逐步进阶而转化为更神圣的：因为身量、智慧和恩宠的增长，记载于圣经，只是为了证明基督确实存在于人性复合体中，从而不给那些主张那里只是一个幻影或人形轮廓、而非真实神圣显现的人留下猜测余地。正因如此，圣经毫无顾忌地记载了祂身上一切人性的偶有特性，甚至吃喝、睡眠、疲倦、养育、身体增长、成长——凡标识人性的一切，除了犯罪倾向。的确，罪是一种流产，而非人性的品质：正如疾病和畸形并非本性原初所固有，而是本性的反常赘生物；同样，趋向罪的活动也应被视为我们内在良善的一种残缺；它本身并非实在之物，我们只在良善缺失时才看见它。因此，那位将我们本性的要素转化为祂神圣能力者，使本性免于残缺和疾病，因为祂自己不允许罪在意志中产生的畸形存在。经上说：\Quote{祂没有犯罪，祂口中也找不到诡诈。}\BibleRef{伯前 2:22}而且，祂身上的这一点不可与任何时间间隔联系起来：因为玛利亚中的那人（智慧在那里建造了她的房屋），虽自然是我们可感复合体的一部分，但随着圣神的降临和至高者能力的庇荫，立刻成为那庇荫能力在本质上所是：因为毫无疑问，较小的受较大的祝福。那么，既然天主性的能力是广大无垠、不可测量的，而人是一个微弱原子，在圣神降临于童贞女、至高者的能力庇荫她的那一刻，由这推动所形成的帐幕并未披上任何人性腐朽之物；而是像起初被构成那样，它一直如此，即使它是人，却仍是圣神、恩宠和能力；我们人性的特殊属性从这丰盛的天主性能力中获得了光辉。\par

人生确实有两个界限：我们起始的那一个，和我们终结的那一个；因此，我们存在的医师有必要在这两端拥抱着我们，不仅抓住终点，也抓住起点，以确保在两者中使受苦者提升。那么，我们在终点一侧所发现发生的事，我们也推论在起点同样发生。正如在终点，藉着降生成人，祂使灵魂虽与身体分离，但已一次与主体（拥有二者者）融合的不可分的天主性，既没有从身体剥落，也没有从灵魂剥落；相反，当它与灵魂同在乐园，并借盗贼之身为全人类开辟了入口时，它藉着身体留在地心，并在那里毁灭了那掌管死亡的——因此祂的\Emph{身体}也因那内在的天主性而被称为\Quote{主}——同样，在起点，我们推论：至高者的能力，藉着圣神的降临（于童贞女）与我们的全部本性融合，既居住在我们的灵魂中——就理性所认为可能居住之处——也与我们的身体结合，使我们的救恩在每个元素上得以完美，而那属于神性的属天无苦，却在作为人的生命的起点和终点都被保存。因此，起点不像我们的起点，终点也不像我们的终点。祂在两者中都显明了祂天主的独立自存；起点没有快乐的玷污，终点也不是因分解的终点。\par

现在，如果我们大声宣讲这一切，并为这一切作证，即基督是天主的德能、天主的智慧，恒常不变、永不朽坏，却进入变化与朽坏之中；祂自己从未被玷污，反而洁净被玷污的；我们犯了什么罪？为什么我们被憎恨？这对抗的新祭台排列是什么意思？我们宣布了另一位耶稣吗？我们暗示了另一位吗？我们提出了别的圣经吗？我们中谁曾胆敢对童贞玛利亚、\OriginalTerm{天主之母}{Mother of God}称\Quote{\OriginalTerm{人之母}{Mother of Man}}？这正是我们听说他们中有些人毫无顾忌地说的。我们编造了三次复活吗？我们应许了千禧年的饕餮吗？我们宣告犹太人的动物祭献将恢复吗？我们再次降低人的希望到地上的耶路撒冷，想象它用更灿烂的石材重建吗？有什么像这样的指控可以加给我们，使我们的团体被视为应当躲避的，并且在某些地方另立一座祭台反对我们，好像我们会玷污他们的圣所？关于此事，我的心曾充满燃烧的义愤；如今我又踏入这座城市，我渴望通过一封信向你们的爱德倾诉我灵魂的苦楚。你们无论圣神引导你们到何处，就留在那里；在天主面前行走；不与血肉商量；不给他们任何人以夸耀的机会，免得他们因你们而夸耀，藉你们生活中的任何事扩大他们的野心。请记住圣父们，你们已获享安息的父亲曾将你们托付给他们，我们蒙天主恩宠被认为配得继承他们；不要挪移我们父辈所立的界石，也不要为了他们更精细的学派而抛弃我们更简朴宣讲的直率。行走在信仰的原始规范中；平安的天主必与你们同在，你们必身心强健。愿天主保守你们不受玷污，这是我们的祈祷。\par

\SourceRecord{Translated by William Moore.；Nicene and Post-Nicene Fathers, Second Series；Vol. 5.；Edited by Philip Schaff and Henry Wace.；Buffalo, NY: Christian Literature Publishing Co.,；1893.；<http://www.newadvent.org/fathers/291117.htm>.}

% structure: p0001 source=h1 target=section reason=page-title

\section{第十八封信}

\Person{圣额我略·尼撒}著\par

致\Person{弗拉维安}\par

天主的人啊，我们这里的情况不妙。某些人无缘无故、毫无来由地对我们产生了恶感，这种恶感的发展已不再是猜测；如今它带着一种只配用于神圣工作的热诚与公开表露出来。而你，一直未曾受到这种烦扰的影响，却过于怠惰，未能扑灭邻人土地上的吞噬火焰；但那些为自己的利益而审慎行事的人，确实会努力制止近处的火灾，通过帮助邻人，确保自己在类似情况下永远不需要帮助。好吧，你会问，我抱怨什么？虔诚已从世界消失；真理已从我们中间逃遁；至于和平，我们过去至少还有这个词在人们嘴边流传；但现在，不仅和平本身不复存在，我们甚至不再保留表达她的词语。但为使你更确切地知道那些激起我们义愤的事情，我将向你简要叙述这整出悲剧。\par

有人曾告诉我，尊敬的\Person{赫拉迪乌斯}对我怀有敌意，并在谈话中对每个人夸耀我给他带来的麻烦。起初我不相信他们的话，只凭我自己以及事情的实际真相来判断。但当每个人都不断向我们讲述同样的话，而事实也证实了他们的报告时，我认为自己有责任不再忽视这种尚无根基与发展趋势的恶感。因此，我写信给你——你的虔诚，以及许多其他能帮助我实现意图的人，激发你们在这件事上的热忱。最后，当我在\Place{塞巴斯蒂亚}结束了纪念最蒙福的\Person{伯多禄}以及与他同时代、人们惯常与他一起纪念的圣殉道者的礼仪后，我正返回自己的主教区，这时有人告诉我\Person{赫拉迪乌斯}本人就在邻近的山丘地区举行殉道者纪念礼仪。起初我继续前行，判断我们的会面更宜在都会本身进行。但当他的一个亲戚不辞辛劳来见我，并向我保证他生病了，我便在得知这消息的地方留下马车；我骑马完成了中间那段路程，那段路如悬崖般陡峭，多石的上坡几乎无法通行。我们必须走完的路程有十五英里。我艰难地行进，时而步行，时而骑马，清晨出发，甚至利用了一部分夜晚，在十二点到一点之间到达了\Place{安杜莫基纳}；因为那就是他与另外两位主教举行会议的地方的名称。从俯瞰这个村庄的山肩上，我们尚在远处便俯视着这次教会的露天集会。我与我的人马缓慢地、步行着、牵着马匹，走过了中间的地面，就在他退回住所时到达了小圣堂。\par

我们立即派信使通知他我们到了；片刻之后，服侍他的执事迎上我们，我们请他即刻告诉\Person{赫拉迪乌斯}，好让我们能与他相处尽可能多的时间，并有机会消除我们之间误解中的一切隔阂。至于我自己，我仍坐在露天之下，等待室内邀请；而我坐在那里，在一个极其不合时宜的时刻，成了所有与会访客的注目对象。时间漫长；睡意袭来，倦怠加剧，旅途劳顿与当天酷热使然；这一切，加上人们盯着我，并向他人指认我，是如此令人痛苦，以致在我身上实现了先知的话：\Quote{我的灵在我里面凄凉。} 我一直保持这种状态直到正午，我由衷地为这次拜访感到后悔，因为给自己带来了这种失礼；而我的反思比敌人对我的伤害更折磨我，它自我争战，转化为对这次冒险的懊悔。最后，通往祭台的道路打开了，我们被允许进入圣所；然而人群被排除在外，尽管我的执事随我进入，用手臂支撑我疲惫的身躯。我向主教阁下致辞，站了一会儿，期望他邀请我就座；但当他没有说出任何这类话时，我便转向一个远处的座位坐下，仍期望他会说出一些友好的话，至少是善意的话；或者至少点头致意。\par

我一度怀抱的任何希望都注定彻底落空。随之而来的是死寂般的沉默，如同悲剧中垂头丧气的面容，困惑、惊愕，以及完全的哑然。那是一段漫长的时间，仿佛在黑夜的黑暗中拖延。我被这种接待如此击垮——他竟不屑于对我施以哪怕那些普通问候中的片言只语，而通常这些问候足以应对偶然相遇的礼节，例如\Quote{欢迎}或\Quote{你从哪里来？}或\Quote{我何幸有此愉悦？}或\Quote{你为何重要事务而来？}——以至于我几乎想把这段沉默的时光画成地下世界的生活图景。不，我谴责这个比喻不恰当。因为在地下世界中，条件完全平等，没有那些在地上生活中引起悲剧的事物扰乱存在。如先知所说，他们的荣耀并不随人而下；每个灵魂，抛弃了大多数人在此世紧紧抓住的东西——他的暴躁、骄傲和自负——以简单无负担的赤裸进入那个低级世界；所以在他们中间找不到今生的任何痛苦。尽管如此，基于这一保留，我当时的处境确实在我看来像一个地下世界，一个阴暗的地牢，一个阴森的刑讯室；尤其是当我想到我们从祖先那里继承了多少社交礼节的宝藏，以及我们将给后代留下多少记载的功绩。其实，我为什么要谈论我们祖先彼此之间的那种亲切态度呢？难怪，既然天生平等，他们不希望彼此占优势，而只想着在谦卑上超过对方。但最打动我心灵的念头是：所有受造物的主，独生子，祂在圣父怀中，祂在起初，祂具有天主的形体，以祂大能的话支撑万有，祂不仅在这点上谦卑自己——在肉体中与人同住——而且甚至欢迎自己的背叛者犹大，当他靠近亲吻祂有福的嘴唇时；当祂进入癞病人西满的家时，祂作为爱众人者，责备主人没有亲吻祂：而我竟不被祂看作与那个癞病人同等；然而我是什么，他又是什么？我无法在我们之间发现任何差别。若从世俗的观点看，他从何处降下，我又躺在何处尘土中？诚然，若必须看待这肉身的生命，这样说也许不会伤害任何人的感情：就我们的血统而论，无论是高贵还是自由，我们的地位大致相当；不过，若寻求真正的自由和高贵，\Emph{即}灵魂的自由和高贵，我们每个人都会被发现有同等的罪的奴仆；每个人同样需要一位除掉他罪过的主；另有一位以祂自己的血将我们从死亡和罪中赎出，救赎了我们，却并不藐视祂所救赎的人，祂将他们从死亡召唤到生命，治愈他们灵魂和肉身的一切软弱。\par

看到这种自负和傲慢的程度如此之大，连高天也几乎显得狭窄（然而我却看不到任何因由或场合引起这种病态心理，如在某些情况下有些人染上它时可以原谅；例如，地位、教育或官职尊严的优越可能使较为虚荣的人膨胀），我无法说服自己保持安静：因为我的内心因整个过程的荒谬而愤愤不平，拒绝接受一切忍耐的理由。那时，若有的话，我对那位神圣宗徒深感钦佩，他如此生动地描绘了我们内在的内战，宣告有一条\BibleQuote{罪的律法在肢体中，反抗理智的律法}，常常使理智成为俘虏，也作它的奴隶。这正是我内心中两种对立情感的排列：一种是因骄傲所致的侮辱而产生的愤怒，另一种则促使平息升起的风暴。靠着天主的恩宠，较坏的倾向未能得胜时，我终于对他说：\Quote{那么，是否我们在此妨碍了你个人舒适所需，该是我们离开的时候了？}当他表示没有身体需要时，我向他说了一些话，尽我所能医治他的恶感。他寥寥数语表示，他对我感到的愤怒是基于我所做的许多伤害。我随即这样回答他：\Quote{谎言在人类中拥有极大的欺骗力量；但在天主的审判中，由此产生的误解将无立足之地。在我与您的关系中，我的良心足以使我盼望，我可能为所有其他罪过获得宽恕，但如果我做过任何损害您的事，但愿这罪永远不被宽恕。}他听了这话很愤怒，不许我再补充我所说之事的证据。\par

此时已过六点，浴室准备妥当，宴席正在摆设，那天正是安息日，且是一位殉道者的纪念日。再看这位福音的门徒是如何效法福音的主：主曾与税吏和罪人一同吃喝，对那些指责祂的人回答说，祂这样做是出于对人的爱；而这位门徒却认为让我们这些经过旅途劳顿、坐在他门口被烈日烤晒、又在他面前遭受他冷淡对待的人与他同席，是一种罪过和玷污。他打发我们拖着疲惫不堪的身躯，原路返回，走了同样的距离；我们几乎在日落时才赶上旅伴，途中还遭遇了许多不幸。因为一阵旋风卷起晴朗天空中的乌云，倾盆大雨把我们淋得湿透；由于天气过于闷热，我们没有防备任何阵雨。然而，靠着天主的恩宠，我们逃脱了，尽管如同从海浪中脱险的沉船水手一样狼狈；我们非常高兴能赶上同伴。\par

我们会合后在那里过夜，终于活着回到了自己的教区；这次会见还带来了一个额外结果，就是他给我们的最后这次侮辱，使之前所有事情的记忆又复活了；你看，我们实在不得不采取措施，为了我们自己，更确切地说，是为了他本人；因为正是由于他以前的企图未被制止，他才肆无忌惮地表现出如此过分的虚荣。因此，我想我们必须做点什么，好让他能改过自新，学会他也不过是人，没有权力侮辱和凌辱那些拥有同样信仰和同等地位的人。试想，假设我们暂时为了论证而承认，我的确做了某些惹恼他的事，那么针对我们是否有过任何审判，来裁决事实或传闻？对于这所谓的伤害，提出了什么证据？引用了哪些教规来控诉我们？有什么合法的主教决议确认了对我们的判决？即使这些程序中的任何一项曾以适当方式进行过，我在教会中的立场固然可能受到威胁，但为了什么教规能允许侮辱一个自由人，并凌辱与他同等地位的人？\Quote{施行公义的审判}，你们这些在这件事上遵循天主法律的人，请说说你们认为我们对所受的凌辱在何处情有可原。如果我们的尊严是根据司铎的管辖权来衡量的，那么大公会议记录的每个司铎的特权都是一样的；更确切地说，因为我们都同样拥有对公教进行监督纠正的职责，所以它是相同的。但如果有人倾向于将我们每个人视为单独的个体，撇开任何司铎的尊严，那么一个人比另一个人在哪方面有优势呢？比如在教育方面，或是在与最高贵、最显赫的世系相连的出身方面，或是在神学方面？这些方面在我身上，要么是相等的，要么至少不逊色。\Quote{但收入呢？}他会说。关于他的收入，我宁愿不被迫提及；只需说这么多就够了：我们自己的收入在一开始是这样，现在也是这样；至于我们收入增长的原因，让别人去探究吧，这收入直到现在还如同被精心培育，几乎每天都在因崇高的事业而增长。那么，他有什么权力侮辱我们呢？他既没有高人一等的出身，也没有超乎众人的地位，更没有雄辩的口才，也从未对我有过恩惠。即使他拥有这一切，轻视出身高贵的人的过错也依然不可原谅。但他没有这些，因此我认为让这种骄傲自大的毛病不加治疗是不合适的；而治疗的方法就是把它压到适当的水平，通过放掉他一点膨胀的自负来缩小其膨胀的尺寸。如何实现这一点，我们交托给天主。\par

\SourceRecord{Translated by William Moore.；Nicene and Post-Nicene Fathers, Second Series；Vol. 5.；Edited by Philip Schaff and Henry Wace.；Buffalo, NY: Christian Literature Publishing Co.,；1893.；<http://www.newadvent.org/fathers/291118.htm>.}
